\section{Pembangunan Model ARIMA}
Setelah data stasioner, dilakukan identifikasi nilai parameter $p$ dan $q$ dengan mengamati pola ACF (untuk $q$) \ref{eq:ma_q_model} dan PACF (untuk $p$) \ref{eq:ar_p_model}, atau dapat pula menggunakan metode otomatis seperti \textit{auto\_arima}. Model kemudian dibangun dan parameternya diestimasi menggunakan \textbf{Maximum Likelihood Estimation (MLE)}. Setelah model dibentuk, dilakukan \textbf{analisis residu} untuk memastikan bahwa kesalahan model bersifat acak \textit{(white noise)}. Hal ini dapat diperiksa melalui plot ACF/PACF dari residu serta uji \textit{Ljung-Box} \ref{eq:ljung_box_q}. Diagnostik ini penting untuk memastikan bahwa model tidak melewatkan pola yang masih ada dalam data. Bila tersedia data uji, model dievaluasi menggunakan metrik prediksi seperti \textbf{RMSE}.